\section{The dynamical model}
\label{sec:model}

\subsection{Variables}

Three quantities evolve together in a dimensionless time $\tau$.

\begin{description}\itemsep3pt
  \item[$\Icap(\tau)$ --- effective capability.] A compressed systems variable
    that rolls together scientific capability, automation, computational
    capacity, coordination and technological leverage. It is not a measure of
    intelligence in the psychometric sense, not a measure of moral worth, and
    not a measure of consciousness. A civilization can grow enormously in
    $\Icap$ without becoming wiser, safer or louder.
  \item[$\Reg(\tau)$ --- regulatory capacity.] Alignment, governance, safety
    culture, institutional learning, epistemic reliability, error correction:
    whatever it takes to stop a runaway. We treat it as a \emph{quality index}
    normalised to $[0,1]$, and that normalisation is what makes it
    commensurable with the growth rates it has to oppose (\S\ref{sec:reg-eq}).
  \item[$\Obs(\tau)$ --- external observability.] Signature strength as defined
    in \S\ref{sec:decomp}, bounded below by $\Ofloor > 0$.
\end{description}

Table~\ref{tab:vars} lists these with their drivers. Figure~\ref{fig:ladder}
maps the usual civilizational stages onto them, and Figure~\ref{fig:schematic}
shows what is coupled to what.

\begin{table}[tb]
\centering\small
\caption{Model variables and drivers. The three states are integrated; the
drivers are exogenous inputs the user sets.}
\label{tab:vars}
\begin{tabular}{@{}llp{7.1cm}@{}}
\toprule
Symbol & Role & Interpretation \\
\midrule
$\Icap$   & state  & effective capability: science, automation, computation, coordination \\
$\Reg$    & state  & regulatory quality index in $[0,1]$: alignment, governance, error correction \\
$\Obs$    & state  & external observability: leakage, waste heat, expansion footprint \\
\midrule
$A$       & driver & ordinary amplification pressure (baseline technological leverage) \\
$\Arec$   & driver & recursive amplification: self-improvement, automated design, compute scaling \\
$\Aref$   & driver & reflective amplification: capability aimed at safety and foresight \\
$Q$       & driver & institutional quality \\
\bottomrule
\end{tabular}
\end{table}

\begin{figure}[tb]
\centering
\begin{tikzpicture}[font=\scriptsize, >={Latex[length=1.8mm]}]
  % Algorithmic hyphenation off: these boxes are too narrow for it to look
  % like anything but a mistake.  Explicit hyphens may still break.
  \hyphenpenalty=10000
  \tikzset{stg/.style={draw, rounded corners=2.5pt, minimum height=7mm,
                       text width=2.1cm, align=center, inner sep=2.2pt,
                       fill=#1}}
  \node[stg=gray!12]                      (s0) {\textbf{0}\\pre-technological};
  \node[stg=blue!10,   right=2mm of s0]   (s1) {\textbf{1}\\early technological};
  \node[stg=blue!17,   right=2mm of s1]   (s2) {\textbf{2}\\planetary};
  \node[stg=orange!20, right=2mm of s2]   (s3) {\textbf{3}\\recursive transition};

  \node[stg=red!13,    above right=3.5mm and 6mm of s3] (s4b)
       {\textbf{4b}\\expansionist\\[1pt]\textit{$\Obs$ stays high}};
  \node[stg=green!17,  right=6mm of s3]                 (s4c)
       {\textbf{4c}\\optimized quiet\\[1pt]\textit{$\Obs$ peaks, then falls}};
  \node[stg=red!24,    below right=3.5mm and 6mm of s3] (s4a)
       {\textbf{4a}\\collapse\\[1pt]\textit{$\Reg$ exhausts}};
  \node[stg=purple!13, right=6mm of s4c]                (s5)
       {\textbf{5}\\post-biological\\[1pt]\textit{poorly constrained}};

  \draw[->] (s0)--(s1);  \draw[->] (s1)--(s2);  \draw[->] (s2)--(s3);
  \draw[->] (s3)--(s4b); \draw[->] (s3)--(s4c); \draw[->] (s3)--(s4a);
  \draw[->] (s4c)--(s5);
\end{tikzpicture}
\caption{Civilizational stages, and the branch at the recursive transition.
Stages 0--2 are pre-recursive and behave the way anyone would expect. Everything
interesting happens at stage~3, where the same machinery can go three different
ways depending on which term wins.}
\label{fig:ladder}
\end{figure}

\begin{figure}[tb]
\centering
\begin{tikzpicture}[font=\small, >={Latex[length=2.2mm]}]
  \tikzset{var/.style={draw, circle, minimum size=13mm, line width=0.7pt, fill=#1},
           drv/.style={font=\footnotesize, align=center, text=black!65}}

  \node[var=blue!16]   (I) at (0,0)      {$\Icap$};
  \node[var=green!18]  (R) at (6.4,0)    {$\Reg$};
  \node[var=orange!22] (O) at (3.2,-3.3) {$\Obs$};

  % exogenous drivers feeding the two upper states
  \node[drv] (dI) at (0,1.85)   {$A,\;\Arec$};
  \node[drv] (dR) at (6.4,1.85) {$\Aref,\;Q$};
  \draw[->, thick, black!55] (dI) -- (I);
  \draw[->, thick, black!55] (dR) -- (R);

  % regulation damps capability
  \draw[->, thick] (R) -- (I);
  \node[anchor=north, font=\footnotesize, yshift=-1.2mm] at ($(I)!0.5!(R)$)
       {damps: $-\,c\Reg\Icap$};

  % the feedback that is deliberately absent -- drawn, then struck out
  \draw[->, dashed, gray!60, line width=0.7pt] (I) to[bend left=26]
        node[pos=0.5, inner sep=0] (nf) {} (R);
  \draw[red!70, line width=1pt] ($(nf)+(-1.7mm,-1.7mm)$) -- ($(nf)+(1.7mm,1.7mm)$);
  \draw[red!70, line width=1pt] ($(nf)+(-1.7mm,1.7mm)$) -- ($(nf)+(1.7mm,-1.7mm)$);
  \node[above=2.4mm of nf, font=\footnotesize, text=gray!45!black]
       {no $\Icap \rightarrow \Reg$ feedback};

  % production and suppression of the signature
  \draw[->, thick] (I) -- node[left, font=\footnotesize, align=right, xshift=-1.5mm]
       {produces\\$+\,\Pi(\Icap)$} (O);
  \draw[->, thick] (R) -- node[right, font=\footnotesize, align=left, xshift=1.5mm]
       {suppresses\\$-\,\Sigma(\Icap,\Reg)$} (O);

  \node[below=6mm of O, font=\footnotesize]
       {$h(\Obs)\,\Psearch \;\longrightarrow\; \Nobs$};
\end{tikzpicture}
\caption{What is coupled to what. Regulation damps capability through
$-c\Reg\Icap$. Observability is produced by capability and suppressed
fractionally by capability \emph{and} regulation acting together --- $\Sigma$
depends on both, even though the arrow has to be drawn from one of them. The
struck-out arc is the coupling this model does \emph{not} have: in
Eq.~\eqref{eq:dR}, regulatory capacity responds to the exogenous drivers and to
erosion by $\Arec$, but never to capability itself. Closing that loop is the
most interesting extension we can see (\S\ref{sec:limitations}).}
\label{fig:schematic}
\end{figure}

\subsection{Capability}

Capability grows two ways: by ordinary accumulation, $a A \Icap$, and by
recursive amplification, $b \Arec F(\Icap)$, where $F$ says how aggressive the
recursion is. It is held back by regulatory friction, $c \Reg \Icap$, and it
eventually runs into physical, energetic and coordination limits, $s_I \Icap^2$:
\begin{equation}
  \dd{\Icap}{\tau} \;=\; a A \Icap \;+\; b \Arec F(\Icap) \;-\; c \Reg \Icap \;-\; s_I \Icap^2 .
  \label{eq:dI}
\end{equation}
Every term is there because some mechanism puts it there. That is the point: the
equation is assembled out of balances you can argue with, not written down
because it integrates nicely.

\subsection{Regulation}
\label{sec:reg-eq}

Regulatory capacity is built by reflective capability and institutional quality,
and eroded by the sheer pace of recursive amplification. The modelling choice
that matters is this: erosion has to be proportional to the regulation that is
actually there, because you cannot erode capacity that does not exist, and
building has to act on whatever headroom is left. So
\begin{equation}
  \dd{\Reg}{\tau} \;=\; \underbrace{(u \Aref + v Q)}_{\textstyle B}\,(1-\Reg)
                 \;-\; \underbrace{(w \Arec + s_R)}_{\textstyle W}\,\Reg .
  \label{eq:dR}
\end{equation}

\begin{proposition}[Regulation is confined to the unit interval]
\label{prop:reg}
If $B \geq 0$ and $W \geq 0$ and $\Reg(0) \in [0,1]$, then $\Reg(\tau) \in [0,1]$
for all $\tau > 0$, and
\begin{equation}
  \Reg_\infty = \frac{B}{B+W} = \frac{u \Aref + v Q}{u \Aref + v Q + w \Arec + s_R} .
  \label{eq:Rinf}
\end{equation}
\end{proposition}

\begin{proof}
At $\Reg = 0$, Eq.~\eqref{eq:dR} gives $\mathrm{d}\Reg/\mathrm{d}\tau = B \geq 0$,
so the lower boundary cannot be crossed downwards. At $\Reg = 1$ it gives
$\mathrm{d}\Reg/\mathrm{d}\tau = -W \leq 0$, so the upper boundary cannot be
crossed upwards. With constant drivers Eq.~\eqref{eq:dR} is linear in $\Reg$,
$\mathrm{d}\Reg/\mathrm{d}\tau = B - (B+W)\Reg$, which gives the stated fixed
point and $\Reg(\tau)=\Reg_\infty+(\Reg_0-\Reg_\infty)e^{-(B+W)\tau}$.
\end{proof}

Two things follow, and both are worth pausing on. First, $\Reg_\infty$ is a
\emph{share}: the fraction of the total pressure on governance that happens to
be constructive. That makes it something you can reason about directly. Second,
because $\Reg$ can never go negative, the damping term $-c\Reg\Icap$ in
Eq.~\eqref{eq:dI} can never flip sign. Let $\Reg$ go negative and failing
governance starts \emph{accelerating} capability, which is not a mechanism
anybody means to model --- it is an arithmetic accident wearing a mechanism's
clothes.

\subsection{Observability}

Observability is produced by energy use, broadcast leakage and expansion
footprint, and suppressed by compression and stealth. The suppression is
fractional: it acts on the signal actually being emitted, above an irreducible
floor.
\begin{equation}
  \dd{\Obs}{\tau} \;=\; \underbrace{p E(\Icap) + q B_{\mathrm c}(\Icap) + r X(\Icap)}_{\textstyle \Pi(\Icap)}
  \;-\; (\Obs - \Ofloor)\,\underbrace{\big(m C(\Icap,\Reg) + n S(\Icap,\Reg)\big)}_{\textstyle \Sigma(\Icap,\Reg)} .
  \label{eq:dO}
\end{equation}

\begin{proposition}[Observability is bounded below by the floor]
\label{prop:obs}
If $\Pi \geq 0$ and $\Sigma \geq 0$ and $\Obs(0) \geq \Ofloor$, then
$\Obs(\tau) \geq \Ofloor$ for all $\tau > 0$. Where $\Sigma > 0$, the
quasi-steady observability is
\begin{equation}
  \Obs^{*} \;=\; \Ofloor + \frac{\Pi(\Icap)}{\Sigma(\Icap,\Reg)} .
  \label{eq:Ostar}
\end{equation}
\end{proposition}

\begin{proof}
At $\Obs = \Ofloor$ the suppression term vanishes identically, leaving
$\mathrm{d}\Obs/\mathrm{d}\tau = \Pi \geq 0$, so the floor cannot be crossed
downwards. Setting $\mathrm{d}\Obs/\mathrm{d}\tau = 0$ in Eq.~\eqref{eq:dO} and
solving for $\Obs$ gives Eq.~\eqref{eq:Ostar}.
\end{proof}

Equation~\eqref{eq:Ostar} is the analytical heart of this paper, and the reason
is that $\Obs^{*}$ is a \emph{ratio}. What it does as capability grows is
decided by how production and suppression scale \emph{against each other},
never by either one alone:
\begin{equation}
  \Pi \sim \Icap, \quad \Sigma \sim \Icap \Reg
    \;\Longrightarrow\; \Obs^{*} \sim \frac{1}{\Reg}
    \qquad\text{(observability independent of capability)},
\end{equation}
\begin{equation}
  \Pi \sim \Icap, \quad \Sigma \sim \Icap^{2} \Reg
    \;\Longrightarrow\; \Obs^{*} \sim \frac{1}{\Icap \Reg}
    \qquad\text{(observability declines with capability)}.
  \label{eq:decline}
\end{equation}
Look at what the second line does. It produces peak-then-decline without anyone
choosing a peaked function --- the shape falls out of the ratio. That is what we
mean by the signature being derived rather than assumed, and \S\ref{sec:results}
shows it happening numerically. \Plausible

\subsection{Interchangeable functional forms}

None of the functional forms above is privileged, and we would rather not
pretend otherwise. Each is drawn from a registry (Table~\ref{tab:registry}), so
asking whether a conclusion survives a change of functions is a menu selection
instead of a rewrite. That is what turns robustness from an aspiration into
something you can actually check in an afternoon.

\begin{table}[tb]
\centering\small
\caption{Interchangeable functional forms. Any combination is admissible; the
solver never refers to a specific choice.}
\label{tab:registry}
\begin{tabular}{@{}lL{3.3cm}L{6.6cm}@{}}
\toprule
Component & Variants & Form \\
\midrule
$F(\Icap)$ recursion & linear, superlinear, Lambert, exponential, saturating &
  $\Icap$; $\Icap^2$; $\Icap e^{\Icap}$; $e^{\Icap}$; $\Icap/(1+\Icap/K)$ \\
$\Obs(\Icap)$ static & A increasing, B decreasing, C peaked, D threshold &
  $1-e^{-\lambda \Icap}$; $e^{-\lambda \Icap}$; $\Icap e^{-\lambda \Icap}$; $\sigma\!\big(k(\Icap-\Icap_c)\big)$ \\
$E$ energy use & linear, Kardashev & $\Icap$; $\Icap^2$ \\
$B_{\mathrm c}$ broadcast & linear, decaying & $\Icap$; $\Icap e^{-\lambda \Icap}$ \\
$X$ expansion & linear, inward & $\Icap$; $0$ \\
$C$, $S$ suppression & linear, superlinear & $\Icap \Reg$; $\Icap^{2}\Reg$ \\
\bottomrule
\end{tabular}
\end{table}

\subsection{Regime classification}

Trajectories are labelled by the criteria in Table~\ref{tab:regimes}, checked in
a fixed order. The order is not commutative --- a trajectory that meets two
criteria gets the earlier label --- and we say so explicitly because it changes
how the phase portrait in \S\ref{sec:results} should be read.

\begin{table}[tb]
\centering\small
\caption{Regime criteria, in evaluation order. Thresholds are model parameters,
not empirical quantities.}
\label{tab:regimes}
\begin{tabular}{@{}llL{6.4cm}l@{}}
\toprule
\# & Regime & Criterion & Strength \\
\midrule
1 & collapse proxy & runaway \emph{and} ($\Reg<\Rmin$ or $\Obs$ peaked then fell below $\Odet$) & \Hypothetical \\
2 & runaway & $\Icap>I_{\mathrm{runaway}}$ or $\Arec/(\Reg+\epsilon)>\Theta$ & \Speculative \\
3 & pre-detectable & $\max\Obs<\Odet$ and $\max\Icap<\Iadv$ & \Plausible \\
4 & optimized low-observable & $\Icap\geq\Iadv$, $\Reg\geq\Rmin$, $\Obs<\Odet$ & \Hypothetical \\
5 & expansionist visible & $\Icap\geq\Iadv$ and $\Obs\geq\Odet$ & \Hypothetical \\
6 & visible technological & $\Pdet>0.01$ & \Plausible \\
7 & uncertain & none of the above & \Plausible \\
\bottomrule
\end{tabular}
\end{table}

Regime~1 is called a \emph{proxy} on purpose. There is no mortality term
anywhere in this model, so what it registers is the conditions under which
collapse should be expected --- not a collapse. We come back to this in
\S\ref{sec:limitations}, because it is a real limitation and not a quibble.
